% Part: incompleteness
% Chapter: introduction
% Section: historical-background

\documentclass[../../../include/open-logic-section]{subfiles}

\begin{document}

\olfileid{inc}{int}{bgr}

\olsection{ঐতিহাসিক পটভূমি}

এই পরিচ্ছেদে আমরা এমন কয়েকটি ঐতিহাসিক বিকাশ সংক্ষেপে আলোচনা করব, যেগুলি
অসম্পূর্ণতা উপপাদ্যগুলিকে যথাযথ প্রসঙ্গে দেখতে সাহায্য করবে। বিশেষত,
গাণিতিক যুক্তিবিদ্যার ইতিহাসের একটি অতি সংক্ষিপ্ত রূপরেখা দেব; তারপর
গণিতের ভিত্তির ইতিহাস সম্বন্ধে কয়েকটি কথা বলব।

\begin{digress}
  ``গাণিতিক যুক্তিবিদ্যা'' কথাটি দ্ব্যর্থক। ``গাণিতিক'' শব্দটি বিষয়বস্তুকে
  বর্ণনা করছে বলে ধরা যায়—যেমন ``গণিতের যুক্তিবিদ্যা'', অর্থাৎ গাণিতিক
  যুক্তিবিচারের নীতিসমূহ; আবার শব্দটি পদ্ধতিকে বর্ণনা করছে বলেও ধরা যায়—যেমন
  ``যুক্তিবিদ্যার গণিত'', অর্থাৎ যুক্তিবিচারের নীতিগুলির গাণিতিক অধ্যয়ন। নিচের
  বিবরণে প্রায়ই একই সঙ্গে উভয় অর্থেই গাণিতিক যুক্তিবিদ্যা আলোচিত হবে।
\end{digress}

যুক্তিবিদ্যার অধ্যয়ন মূলত শুরু হয় অ্যারিস্টটলের হাতে; তিনি আনুমানিক
৩৮৪--৩২২ \textsc{bce}-তে জীবিত ছিলেন। তাঁর \emph{Categories}, \emph{Prior
  analytics} এবং \emph{Posterior analytics} গ্রন্থে বৈজ্ঞানিক যুক্তিবিচারের
নীতিগুলির সুসংবদ্ধ আলোচনা আছে; বিশেষ করে ন্যায়ানুমানের একটি বিস্তৃত ও
সুশৃঙ্খল অধ্যয়ন আছে।

মধ্যযুগ জুড়ে পাণ্ডিত্যবাদী দর্শনে অ্যারিস্টটলের যুক্তিবিদ্যাই প্রাধান্য
পেয়েছিল; এমনকি অষ্টাদশ শতাব্দীতেও কান্ট বলেছিলেন, অ্যারিস্টটলের যুক্তিবিদ্যা
নিখুঁত এবং তার কোনো সংশোধন দরকার নেই। কিন্তু গাণিতিক যুক্তিবিচারের একেবারে
উপরিতলের দিকগুলি ছাড়া অন্য কিছু মডেল করার পক্ষে ন্যায়ানুমানের তত্ত্ব অত্যন্ত
সীমিত। তার এক শতাব্দী আগে নিউটনের সমসাময়িক লাইবনিৎস যুক্তিবিচারের জন্য একটি
সম্পূর্ণ ``কলন'' কল্পনা করেছিলেন এবং এমন কলন নির্মাণের কয়েকটি প্রাথমিক পদক্ষেপও
নিয়েছিলেন; মূলত তিনি বচনমূলক যুক্তিবিদ্যার একটি রূপ বর্ণনা করেছিলেন।

উনিশ শতক যুক্তিবিদ্যার ইতিহাসে এক সন্ধিক্ষণ। ১৮৫৪ সালে জর্জ বুল
\emph{The Laws of Thought} লেখেন; এতে বচনমূলক যুক্তিবিদ্যার এমন এক বিস্তৃত
বীজগাণিতিক অধ্যয়ন ছিল যা আধুনিক উপস্থাপনা থেকে খুব দূরে নয়। ১৮৭৯ সালে
গটলব ফ্রেগে তাঁর \emph{Begriffsschrift} (ধারণালিপি) প্রকাশ করেন; এটি
বচনমূলক যুক্তিবিদ্যাকে পরিমাণসূচক ও সম্বন্ধ দিয়ে প্রসারিত করে এবং তাই
প্রথম-ক্রমের যুক্তিবিদ্যাকে অন্তর্ভুক্ত করে। বস্তুত, ফ্রেগের যৌক্তিক ব্যবস্থায়
উচ্চতর-ক্রমের যুক্তিবিদ্যা এবং আরও অনেক কিছুও ছিল। তাঁর \emph{Basic Laws
  of Arithmetic}-এ ফ্রেগে দেখাতে চেয়েছিলেন যে কেবল যৌক্তিক প্রকল্প থেকে তাঁর
Begriffsschrift-এ সমগ্র পাটিগণিত নিষ্পাদন করা যায়। দুর্ভাগ্যক্রমে, রাসেল
১৯০২ সালে দেখান যে এই প্রকল্পগুলি অসঙ্গত। কিন্তু অসঙ্গত স্বতঃসিদ্ধটি বাদ দিলে
ফ্রেগে প্রায় একক প্রচেষ্টায় আধুনিক যুক্তিবিদ্যা উদ্ভাবন করেছিলেন—এ এক বিস্ময়কর
কীর্তি। বুলের পর পিয়ার্স ও শ্র্যোডার-সহ বীজগাণিতিক মনোভাবাপন্ন চিন্তকেরাও
স্বাধীনভাবে পরিমাণসূচকীয় যুক্তিবিদ্যা বিকশিত করেছিলেন।

এবার গণিতের ভিত্তির বিকাশে আসা যাক। যুক্তিবিদ্যা গণিতে গুরুত্বপূর্ণ ভূমিকা
নেয় বলে এই বিকাশের সঙ্গে সদ্য-বর্ণিত ঘটনাগুলির যথেষ্ট আন্তঃক্রিয়া ছিল। যেমন,
ফ্রেগে তাঁর যুক্তিবিদ্যা নির্মাণ করেছিলেন স্পষ্টতই দেখানোর উদ্দেশ্যে যে সমগ্র
গণিত কেবল তাঁর যৌক্তিক কাঠামোর উপর প্রতিষ্ঠিত হতে পারে; বিশেষ করে তিনি দেখাতে
চেয়েছিলেন, কান্টের দাবির মতো গণিতের সত্যগুলি অবধারণাপূর্ব emph{সংশ্লেষক}
নয়, বরং অবধারণাপূর্ব emph{বিশ্লেষক}।

অনেকে মনে করেন, যথার্থ গণিতের জন্ম হয়েছিল গ্রিকদের হাতে। আনুমানিক ৩০০
B.C.-তে লেখা ইউক্লিডের \emph{Elements} কঠোরতা ও সূক্ষ্মতার উপর গুরুত্বসহ
গ্রিক গণিতের এক পরিণত নিদর্শন। ইউক্লিডের \emph{Elements}-এর সংজ্ঞা ও
প্রমাণগুলি আজও উচ্চবিদ্যালয়ের জ্যামিতি পাঠ্যবইয়ে প্রায় অবিকৃত অবস্থায় টিকে
আছে (যে পরিমাণে উচ্চবিদ্যালয়ে এখনও জ্যামিতি পড়ানো হয়)। গাণিতিক যুক্তিবিচারের
এই আদর্শকে শুধু গণিতে নয়, দর্শনের বিভিন্ন শাখাতেও কঠোর যুক্তির দৃষ্টান্ত বলে
মানা হয়েছে। (স্পিনোজা এমনকি নৈতিক ও ধর্মীয় যুক্তিও ইউক্লিডীয় রীতিতে
উপস্থাপন করেছিলেন—দেখতে বেশ অদ্ভুত!)

সপ্তদশ শতাব্দীতে নিউটন ও লাইবনিৎস কলন আবিষ্কার করেন। (অগ্রাধিকারের প্রশ্নে
শতাব্দীর পর শতাব্দী তীব্র বিবাদ চলেছে, কিন্তু আজ অধিকাংশ গবেষক মনে করেন,
দুটি বিকাশ প্রধানত স্বাধীন ছিল।) কলনে, যেমন, অসীম ক্ষুদ্র রাশির অসীম যোগফল
নিয়ে যুক্তিবিচার করা হয়; এই বৈশিষ্ট্যগুলি বিশপ বার্কলির সমালোচনাকে উসকে দেয়।
তিনি যুক্তি দিয়েছিলেন, ঈশ্বরে বিশ্বাস তাঁর সময়ের গণিতের চেয়ে কম যুক্তিসঙ্গত
নয়। অষ্টাদশ শতাব্দীতে কলনের পদ্ধতি ব্যাপকভাবে ব্যবহৃত হয়; যেমন লেওনার্ড
অয়লার অসীম যোগফল-যুক্ত গণনা করে নাটকীয় ফল পেয়েছিলেন।

উনিশ শতকে গণিতবিদেরা কলনকে আরও দৃঢ় ভিত্তির উপর রেখে বার্কলির সমালোচনার
জবাব দেওয়ার চেষ্টা করেন। কোশি, ভাইয়ারস্ট্রাস, বলৎসানো প্রমুখের প্রচেষ্টা থেকে
``এপসিলন ও ডেল্টা''-র সাহায্যে সীমা, অবিচ্ছিন্নতা, অন্তরীকরণ ও সমাকলনের
আধুনিক সংজ্ঞাগুলি আসে—অর্থাৎ অতিক্ষুদ্র রাশির কোনো উল্লেখ ছাড়াই। শতাব্দীর
পরের দিকে গণিতবিদেরা আরও এগিয়ে বাস্তব সংখ্যা-সহ কলনের সব দিককে স্বাভাবিক
সংখ্যার সাহায্যে ব্যাখ্যা করতে চাইলেন। (ক্রোনেকার: ``ঈশ্বর পূর্ণসংখ্যা সৃষ্টি
করেছেন; বাকি সব মানুষের কাজ।'') ১৮৭২ সালে ডেডেকিন্ড ``Continuity and the
irrational numbers'' লেখেন; সেখানে তিনি দেখান কীভাবে মূলদ সংখ্যার সেট হিসেবে
বাস্তব সংখ্যা ``নির্মাণ'' করা যায় (মূলদ সংখ্যাকে, তোমরা জানো, স্বাভাবিক সংখ্যার
জোড়া হিসেবে দেখা যায়)। ১৮৮৮ সালে তিনি লেখেন ``Was sind und was sollen die
Zahlen'' (মোটামুটি, ``স্বাভাবিক সংখ্যা কী, এবং তাদের কী হওয়া উচিত?''); এর লক্ষ্য
ছিল বিশুদ্ধ ``যৌক্তিক'' ভাষায় স্বাভাবিক সংখ্যা ব্যাখ্যা করা। ১৮৮৭ সালে ক্রোনেকার
``\"Uber den Zahlbegriff'' (``সংখ্যার ধারণা সম্বন্ধে'') লেখেন, যেখানে তিনি
পূর্ণসংখ্যার সাহায্যে সব গাণিতিক বস্তু উপস্থাপনের কথা বলেন; ১৮৮৯ সালে জুসেপ্পে
পেয়ানো স্বাভাবিক সংখ্যার জন্য বিধিবদ্ধ প্রতীকী স্বতঃসিদ্ধ দেন।

উনিশ শতকের শেষভাগে অসীম নিয়ে কাজ করার ক্ষেত্রেও নতুন সাহস দেখা দেয়। তার
আগে অসীমধর্মী বস্তু ও গঠন (যেমন স্বাভাবিক সংখ্যার সেট) খুব সাবধানে ব্যবহার করা
হতো; ``অসীমসংখ্যক'' বলতে বোঝানো হতো ``যতগুলি চাই'', আর ``সীমায় এগিয়ে যায়''
বলতে বোঝানো হতো ``যত কাছে চাই তত কাছে যায়''। কিন্তু গেয়র্গ কান্টর দেখান যে
অসীমকে সরাসরি একটি সম্পূর্ণ সত্তা হিসেবে গ্রহণ করা সম্ভব। কান্টর, ডেডেকিন্ড
প্রমুখের কাজ গণিতের সেই সাধারণ সেটতাত্ত্বিক ধারণা প্রবর্তনে সাহায্য করে, যা এখন
ব্যাপকভাবে গৃহীত।

এবার বিংশ শতাব্দীর যুক্তিবিদ্যা ও ভিত্তিবিষয়ক বিকাশে আসি। ১৯০২ সালে রাসেল
ফ্রেগের যৌক্তিক ব্যবস্থায় কূটাভাসটি আবিষ্কার করেন। ১৯০৪ সালে জার্মেলো তথাকথিত
``নির্বাচন স্বতঃসিদ্ধ'' ব্যবহার করে কান্টরের সুবিন্যাস নীতি প্রমাণ করেন; এই
স্বতঃসিদ্ধের বৈধতা নিয়ে বিস্তর বিতর্ক হয়। ১৯১০ থেকে ১৯১৩ সালের মধ্যে রাসেল
ও হোয়াইটহেডের \emph{Principia Mathematica}-র তিনটি খণ্ড প্রকাশিত হয় এবং
যৌক্তিক ভিত্তির উপর গণিত প্রতিষ্ঠার ফ্রেগীয় কর্মসূচিকে প্রসারিত করে। কিন্তু
রাসেল ও হোয়াইটহেডকে এমন দুটি নীতি গ্রহণ করতে হয়েছিল যেগুলিকে বিশুদ্ধ যৌক্তিক
বলে সমর্থন করা কঠিন: অসীমতার একটি স্বতঃসিদ্ধ এবং ``হ্রাসযোগ্যতা''-র একটি
স্বতঃসিদ্ধ। ১৯০০-এর দশকে পঁয়কারে গণিতে ``অপ্রেডিকেটিভ সংজ্ঞা'' ব্যবহারের
সমালোচনা করেন; ১৯১০-এর দশকে ব্রাউয়ার বর্জিত মধ্যমের নিয়ম
($!A \lor \lnot !A$) পরিহার করে ``স্বজ্ঞাবাদী'' ভিত্তিতে সমগ্র গণিত নতুন করে
প্রতিষ্ঠার প্রস্তাব দিতে শুরু করেন।

সত্যিই বিচিত্র সময়!{} সমগ্র গণিতকে যুক্তিবিদ্যায় হ্রাস করার কর্মসূচি এখন
``যুক্তিবাদ'' নামে পরিচিত এবং উপরোক্ত সমস্যাগুলির জন্য সাধারণত ব্যর্থ বলে গণ্য।
স্বজ্ঞাবাদী মানসিক নির্মাণের সাহায্যে গণিত বিকাশের কর্মসূচিকে ``স্বজ্ঞাবাদ'' বলা
হয় এবং সাধারণ গাণিতিক চর্চার উপর তা অতিরিক্ত কঠোর বাধা আরোপ করে বলে মনে করা
হয়। শতাব্দীর সন্ধিক্ষণে সর্বকালের সবচেয়ে প্রভাবশালী গণিতবিদদের অন্যতম ডাভিড
হিলবার্ট কান্টর ও ডেডেকিন্ড-প্রবর্তিত নতুন বিমূর্ত পদ্ধতির প্রবল সমর্থক ছিলেন:
``কান্টর আমাদের জন্য যে স্বর্গ সৃষ্টি করেছেন, সেখান থেকে কেউ আমাদের তাড়াতে
পারবে না।'' একই সঙ্গে তিনি এই নতুন পদ্ধতিগুলির (বিচিত্রভাবে, এখন যেগুলিকে
``ধ্রুপদি'' বলা হয়) ভিত্তিবিষয়ক সমালোচনা সম্বন্ধে সচেতন ছিলেন। দুই দিক বজায়
রাখার একটি উপায় তিনি প্রস্তাব করেন:
\begin{enumerate}
\item বিধিবদ্ধ স্বতঃসিদ্ধ ও নিয়ম দিয়ে ধ্রুপদি পদ্ধতিগুলিকে উপস্থাপন করো;
  একটি স্বতঃসিদ্ধমূলক ব্যবস্থায় গাণিতিক প্রশ্নগুলিকে !!{formula}s হিসেবে
  উপস্থাপন করো।
\item নিরাপদ, ``সসীমধর্মী'' পদ্ধতি দিয়ে প্রমাণ করো যে এই বিধিবদ্ধ
  অবরোহী ব্যবস্থাগুলি সঙ্গত।
\end{enumerate}

হিলবার্টের কাজ প্রথম লক্ষ্য পূরণের পথে অনেক দূর এগিয়েছিল। ১৮৯৯ সালে তাঁর
বিখ্যাত \emph{Foundations of geometry} গ্রন্থে তিনি জ্যামিতির জন্য এই কাজ
করেছিলেন। পরের বছরগুলিতে তিনি এবং তাঁর বহু ছাত্র ও সহযোগী গণিতের অন্যান্য
ক্ষেত্রে জ্যামিতির কাজটির অনুরূপ কাজ করেন। হিলবার্ট নিজে পাটিগণিত ও বিশ্লেষণের
জন্য স্বতঃসিদ্ধ-ব্যবস্থা দেন। জার্মেলো সেটতত্ত্বের একটি স্বতঃসিদ্ধায়ন দেন;
ফ্রেঙ্কেল, স্কোলেম, ফন নয়মান প্রমুখ পরে সেটিকে প্রসারিত করেন। ১৯২০-এর দশকের
মাঝামাঝি ``সমগ্র'' গণিতের স্বতঃসিদ্ধায়ন বলে দাবি করতে পারে এমন দুটি পদ্ধতি
ছিল—রাসেল ও হোয়াইটহেডের \emph{Principia mathematica} এবং জার্মেলো--ফ্রেঙ্কেল
সেটতত্ত্ব নামে পরিচিত ব্যবস্থাটি।

১৯২১ সালে হিলবার্ট এই ব্যবস্থাগুলির সঙ্গতি প্রমাণের লক্ষ্য নিয়ে একটি গবেষণা
কর্মসূচি শুরু করেন। তাঁর কয়েকজন ছাত্র—বিশেষত বের্নাইস, আকারমান এবং পরে
গেন্ৎসেন—এই কর্মসূচিতে তাঁকে সাহায্য করেন। মূল ধারণাটি ছিল, গণিতে অসঙ্গতির
একটি !!{derivation} সম্ভব কি না সেই প্রশ্নকে সম্ভাব্য প্রতীকক্রমের একটি
সমাবেশগত সমস্যা হিসেবে প্রকাশ করা। অর্থাৎ পাটিগণিত, বিশ্লেষণ বা সেটতত্ত্বের
কোনো স্বতঃসিদ্ধ-ব্যবস্থার স্বতঃসিদ্ধগুলি থেকে, ধরা যাক, $!A \land
\lnot !A$-এর সঠিক !!{derivation} হওয়ার মাপকাঠি পূরণ করে এমন সম্ভাব্য
বাক্যক্রম নিয়ে প্রশ্নটি গঠিত হবে। এমন প্রতীকক্রমের অসম্ভবতার প্রমাণটি নিজেও
একটি গাণিতিক প্রমাণ বলে এই স্বতঃসিদ্ধমূলক ব্যবস্থাগুলিতে বিধিবদ্ধ করা সম্ভব
হবে। অন্যভাবে বললে, এমন একটি বাক্য~$\OCon$ থাকবে যা, ধরা যাক, পাটিগণিতের
সঙ্গতি ব্যক্ত করে। তদুপরি, বিশেষত যদি তার প্রমাণে কেবল অত্যন্ত সীমিত
``সসীমধর্মী'' উপায় লাগে, তবে বাক্যটি সংশ্লিষ্ট ব্যবস্থাতেই প্রমাণযোগ্য হওয়া
উচিত।

আরেকটি লক্ষ্য—নির্মিত স্বতঃসিদ্ধ-ব্যবস্থাগুলি যেন প্রতিটি গাণিতিক প্রশ্নের
মীমাংসা করে—দুভাবে নির্দিষ্ট করা যায়। প্রথমভাবে বলা যায়: গণিতের কোনো
স্বতঃসিদ্ধ-ব্যবস্থার ভাষায় প্রতিটি বাক্য~$!A$-এর ক্ষেত্রে স্বতঃসিদ্ধগুলি থেকে
$!A$ অথবা $\lnot !A$ প্রমাণযোগ্য। তা সত্য হলে স্বতঃসিদ্ধের ভিত্তিতে প্রমাণও
খণ্ডনও করা যায় না এমন কোনো বাক্য থাকবে না, স্বতঃসিদ্ধগুলি মীমাংসা করে না এমন
কোনো প্রশ্ন থাকবে না। এই ধর্মবিশিষ্ট স্বতঃসিদ্ধ-ব্যবস্থাকে
\emph{পূর্ণ} বলা হয়। অবশ্য কোনো নির্দিষ্ট বাক্যের ক্ষেত্রে দুই বিকল্পের কোনটি
সত্য তা নির্ধারণ করা তবুও কঠিন হতে পারে। কিন্তু নীতিগতভাবে তা করার একটি পদ্ধতি
থাকা উচিত। বস্তুত হিলবার্ট যে স্বতঃসিদ্ধ ও !!{derivation} ব্যবস্থাগুলি
বিবেচনা করেছিলেন, সেগুলির ক্ষেত্রে পূর্ণতা এমন পদ্ধতির অস্তিত্ব থেকে অনুসৃত
হতো—যদিও হিলবার্ট তা বুঝতে পারেননি। প্রশ্নটির দ্বিতীয়, আরও প্রবল পাঠ হলো:
এমন একটি যান্ত্রিক গণনাপদ্ধতি থাকবে যা কোনো বাক্য~$!A$ দেওয়া হলে সেটি
স্বতঃসিদ্ধগুলি থেকে নিষ্পাদনযোগ্য কি না নির্ধারণ করবে।
% BN-SRC-198 TARGET-CORRECTION-BEGIN
\par\smallskip\noindent\textnormal{\small\textbf{উৎস-সংশোধন (BN-SRC-198):} হিমায়িত অনুচ্ছেদটি একে ``দ্বিতীয় লক্ষ্য'' বলে, কিন্তু ঠিক আগের সংখ্যায়িত তালিকায় দ্বিতীয় লক্ষ্য হলো বিধিবদ্ধ ব্যবস্থার সঙ্গতি প্রমাণ করা। বাংলা পাঠে এই পৃথক পূর্ণতা/নির্ণেয়তার উদ্দেশ্যটিকে ``আরেকটি লক্ষ্য'' বলা হয়েছে।} % BN-SRC-198 TARGET-CORRECTION-END

১৯৩১ সালে গ্যোডেল দুটি ``অসম্পূর্ণতা উপপাদ্য'' প্রমাণ করে দেখান যে এই
কর্মসূচির সব লক্ষ্য সফল হতে পারে না। যথেষ্ট শক্তিশালী, সঙ্গত ও কার্যকরভাবে
স্বতঃসিদ্ধায়িত গাণিতিক ব্যবস্থা পূর্ণ হতে পারে না; বিশেষত দ্বিতীয়
অসম্পূর্ণতা উপপাদ্যের শর্তগুলি পূরণকারী এমন ব্যবস্থা নিজের সঙ্গতি-বাক্য প্রমাণ
করতে পারে না।
% BN-SRC-199 TARGET-CORRECTION-BEGIN
\par\smallskip\noindent\textnormal{\small\textbf{উৎস-সংশোধন (BN-SRC-199):} হিমায়িত সারাংশটি কোনো শর্ত ছাড়াই বলে যে গণিতের কোনো স্বতঃসিদ্ধ-ব্যবস্থা পূর্ণ নয়, তারপর নিজের সঙ্গতি-বাক্য প্রমাণ বা খণ্ডন কোনোটিই করা যায় না বলে। অসম্পূর্ণতা উপপাদ্যগুলির জন্য কার্যকর স্বতঃসিদ্ধায়ন, যথেষ্ট গাণিতিক শক্তি ও উপযুক্ত সঙ্গতি-অনুমান দরকার; দ্বিতীয় উপপাদ্যের সাধারণ সিদ্ধান্ত হলো নিজের সঙ্গতি প্রমাণ করা যায় না, খণ্ডনও করা যায় না—এটি নিছক সঙ্গতি থেকে অনুসৃত নয়। বাংলা সারাংশে অধ্যায়ে পরে নির্দিষ্ট করা সঠিক শর্ত ও নিশ্চিত সিদ্ধান্ত রাখা হয়েছে।} % BN-SRC-199 TARGET-CORRECTION-END

এটি হিলবার্টের মূল কর্মসূচির উপর মারণ আঘাত হেনেছিল। কিন্তু গণিতে প্রায়ই যেমন
হয়, তেমনই এটি গবেষণার রোমাঞ্চকর নতুন পথও খুলে দেয়। যদি সমগ্র গণিতকে ধারণ
করা একটিমাত্র বিধিবদ্ধ ব্যবস্থা না থাকে, তবে আরও সীমিত ব্যবস্থা নির্মাণ করে
তাতে কী প্রমাণ করা যায় তা অনুসন্ধান করা সার্থক। এই ব্যবস্থাগুলির সঙ্গতি
প্রতিষ্ঠার জন্য কম সীমিত প্রমাণপদ্ধতি নির্মাণ করা এবং তাদের সঙ্গতি প্রমাণ করা
কত কঠিন তা মাপার উপায় খোঁজাও সার্থক। গ্যোডেল দেখিয়েছেন যে (প্রায়) প্রতিটি
বিধিবদ্ধ ব্যবস্থারই এমন প্রশ্ন আছে যার মীমাংসা সে করতে পারে না; তাই কোনো
নির্দিষ্ট ব্যবস্থা মীমাংসা করতে পারে না এমন ``আকর্ষণীয়'' প্রশ্ন খোঁজা এবং
সেগুলি মীমাংসার জন্য ব্যবস্থা কতটা শক্তিশালী হওয়া চাই তা নির্ণয় করা যুক্তিসঙ্গত।
আজও যুক্তিবিদেরা প্রমাণতত্ত্ব নামের নতুন গাণিতিক শাস্ত্রে এই প্রশ্নগুলির অনুসন্ধান
করে চলেছেন।

\end{document}
